% !TeX program = xelatex
% Copyright (C) 2026 bsxiaocai
% SPDX-License-Identifier: LPPL-1.3c
% Example under LPPL 1.3c. See LICENSE and NOTICE.md.
% Compile from the repository root.
% 逐页操作说明：docs/templates/academic.md；全部尺寸：docs/customization.md。
\documentclass[aspectratio=169]{ctexbeamer}
\usetheme{HUAS}
\HUASsetup{style=academic,logo=true,sectionpage=true,faculty={学术交流}}
\input{examples/local-assets.tex}
% 微调放在导言区；先取消下面注释，再只改需要的参数。
% \HUASsetlayout{frametitle padding=2mm,hero rule width=3cm}
% \setbeamerfont{frametitle}{size=\Large} % 正文标题字号；加大后减少正文可用高度。
% 主标题、副标题、作者、机构、日期分别在以下五行替换。
\title{从阅读问题到研究证据}
\subtitle{学术汇报示例}
\author{汇报者}
\institute{湖南文理学院}
\date{\today}
\begin{document}
\HUASmaketitle
\begin{frame}{汇报结构}
  % 修改 section 后编译两遍；删 hideallsubsections 可显示已定义的小节。
  \tableofcontents[hideallsubsections]
\end{frame}
\section{问题与方法}
\begin{frame}{研究问题：学生如何解释文本？}
  \begin{block}{研究对象}
    观察学生从选择语句到形成解释的过程。
  \end{block}
  \begin{itemize}
    \item 问题：解释是否有明确的文本依据？
    \item 材料：阅读批注、讨论记录与修订稿。
    \item 边界：本示例展示汇报结构，不提供实证结果。
  \end{itemize}
\end{frame}
\begin{frame}{方法：建立可复核的分析过程}
  \begin{enumerate}
    \item 保留原始记录，明确每段材料的来源。
    \item 分别标记观点、证据与推理过程。
    \item 比较不同判断，并记录分歧的处理方式。
  \end{enumerate}
  \medskip
  用比例说明证据覆盖范围：
  % 显示公式使用原生数学环境；supported/total 是直立文字下标。
  \[
    r=\frac{n_{\mathrm{supported}}}{n_{\mathrm{total}}}
  \]
  其中分子为有文本依据的解释数，分母为全部解释数。
\end{frame}
\section{分析与讨论}
\begin{frame}{分析框架：让判断与材料对应}
  \begin{table}
    \caption{观察维度与材料位置}
    % 1.5 为行高倍数；lll 表示三列左对齐，长文字可改 p{4cm} 固定列宽。
    \renewcommand{\arraystretch}{1.5}
    \begin{tabular}{lll}
      \hline
      \textbf{维度} & \textbf{观察问题} & \textbf{材料} \\
      \hline
      观点 & 判断是否清楚？ & 讨论记录 \\
      证据 & 语句是否相关？ & 阅读批注 \\
      推理 & 如何联系观点？ & 修订稿 \\
      \hline
    \end{tabular}
  \end{table}
  \begin{alertblock}{报告时需要说明}
    分析框架与研究发现应分别呈现；未验证的判断不要写成结论。
  \end{alertblock}
\end{frame}
% 此命令自己生成一页，必须在 frame 外；第二参数中的 \\ 强制换行。
\HUASkeypoint{讨论与结论}{结论回到证据，\\解释保留边界。}
\begin{frame}{参考资料}
  \begin{thebibliography}{9}
    % 每条 bibitem 的键应唯一；正文可用 \cite{guide}，添加引用后编译两遍。
    \bibitem{guide}Till Tantau, Joseph Wright, Vedran Miletić.
      \newblock The beamer class: User Guide.
      \newblock TeX Live 随附文档。
    \bibitem{project}HUAS-Beamer 项目.
      \newblock 主题使用说明与阶段记录.
      \newblock 本仓库 README 与 docs。
  \end{thebibliography}
\end{frame}
\HUASclosingframe
\end{document}
